\documentclass[10pt,a4paper]{amsart}
\usepackage[margin=19mm]{geometry}
\usepackage{amsmath,amssymb,mathtools}
\usepackage[scr=rsfs,cal=euler]{mathalfa}
\usepackage{xfrac}
\usepackage[unicode,hidelinks,bookmarks=false]{hyperref}
\newcommand{\ie}{i.e.\ }
\newcommand{\cf}{cf.\ }
\newcommand{\resp}{respectively\ }
\newcommand{\Subsection}{\subsection}
\providecommand{\nobreakdash}{\nobreak\hbox{-}}
\setlength{\emergencystretch}{2em}
\multlinegap0pt
\usepackage{bm}
\usepackage{bbm}
\usepackage{dsfont}
\usepackage{xspace}
\usepackage{cancel}
\usepackage{mathtools}
\usepackage{amsthm}
\makeatletter
\@ifundefined{theorem}{\newtheorem{theorem}{Theorem}}{}
\@ifundefined{lemma}{\newtheorem{lemma}[theorem]{Lemma}}{}
\makeatother
\title[Some convolution identities for mock modular forms arising f]{Some convolution identities for mock modular forms arising from the theory of holomorphic projection}
\author{Jonathan G. Bradley-Thrush, Frank Garvan, Jayashree Kalita,, Larry Rolen}
\date{Source: arXiv:2608.13324v1 (CC BY 4.0)}
\begin{document}
\maketitle

\noindent\emph{Source and licence.} Some convolution identities for mock modular forms arising from the theory of holomorphic projection. Version arXiv:2608.13324v1 (CC BY 4.0), distributed under the Creative Commons Attribution 4.0 International licence, as recorded in the arXiv licence metadata for this exact version and shown on the abstract page https://arxiv.org/abs/2608.13324v1. Full rights notice: licenses/arxiv-2608.13324-CC-BY-4.0.txt.\par\smallskip

\noindent\textit{Editorial scope.} Counted unit: the Lemma whose source label is \texttt{lem:double\_ser\_trans} in the article (source0.tex lines 1732--1738, source label \texttt{lem:double\_ser\_trans}), delivered together with the article's own complete original proof of it (source0.tex lines 1740--1788, one \texttt{proof} environment of 49 lines). No mathematics is written, repaired or re-derived: statement and proof are the article's own text line for line. Required context retained uncounted: none. Every definition, display and label of those lines is retained unchanged. Omitted from this package and not redistributed: 1--1731, 1739--1739, 1789--2732. Nothing inside a retained excerpt is omitted or altered: every retained body is delivered exactly as the article's lines are written, so no omission receipt of any kind accompanies this package. Citations: the retained excerpts carry no citation command at all, so no bibliography entry of the article is transcribed and no reference is redistributed. Reference adaptation: none was needed and none was made; every reference inside the delivered unit resolves inside it, so no unresolved reference or citation is produced. Printed slips kept literal: no printed word, symbol, hypothesis or number of the counted statement or of its proof was altered. Layout and class: the article's own class is \texttt{amsart} with packages including amsmath, mathtools, amsthm, amssymb, verbatim, xcolor, hyperref, enumitem, graphicx, fullpage, todonotes, tikz, fancyhdr, color; the wrapper is the lane's pinned \texttt{amsart} editorial wrapper at 10pt on A4, which restates the theorem-like environments the retained text needs and adds the article's own macro definitions that the retained text needs (none), with extra wrapper packages (bm, bbm, dsfont, xspace, cancel, mathtools). The delivered numbering is the unit's own. Assets: no retained span contains a figure environment, a graphics-inclusion command, a PDF-page inclusion, a table or a TikZ picture; no image file is copied into this package, the package declares no asset dependency and no image file is redistributed with it. Licence: version arXiv:2608.13324v1 of this article is distributed under the Creative Commons Attribution 4.0 International licence, as recorded in the arXiv licence metadata for this exact version and shown on the abstract page https://arxiv.org/abs/2608.13324v1. A full rights notice is delivered at licenses/arxiv-2608.13324-CC-BY-4.0.txt. Characters: every retained excerpt is pure ASCII, so the delivered engine, XeLaTeX with its default Latin Modern text font, reports no missing character.
\begin{lemma} \label{lem:double_ser_trans}
For any function $\varphi:\mathbb{Z}^2 \rightarrow \mathbb{C}$ such that the series below converge absolutely, the following identities hold:
\begin{align}
\sum_{m=1}^\infty \sum_{n=1}^{\lfloor {12 \over 7}m \rfloor} \left({-4 \over m}\right) \left({n \over 6}\right) \varphi(m,n) &= \sum_{m=1}^\infty \sum_{n=1}^{\lfloor {3 \over 2}m \rfloor} \left({-4 \over m}\right) \left({n \over 6}\right) \Big( \varphi(m,n) - \varphi(7m-4n,12m-7n) \Big), \label{eqn:double_ser_trans1} \\
\sum_{m=1}^\infty \sum_{n=1}^{\lfloor {24 \over 7}m \rfloor} \left({-4 \over m}\right) \left({n \over 6}\right) \varphi(m,n) &= \sum_{m=1}^\infty \sum_{n=1}^{3m} \left({-4 \over m}\right) \left({n \over 6}\right) \Big( \varphi(m,n) - \varphi(7m-2n,24m-7n) \Big). \label{eqn:double_ser_trans2}
\end{align}
\end{lemma}

\begin{proof}
Let $S$ denote the double series on the left-hand side of~\eqref{eqn:double_ser_trans1}. 
The unimodular substitution
\begin{equation} \label{eqn:unimodular_sub}
(m,n) \mapsto (2m-n, 3m-2n)
\end{equation}
changes the inequalities
\[ n \geq 0 \qquad \text{and} \qquad n \leq {7m \over 12}, \]
which define the range of the summation on the left-hand side of~\eqref{eqn:double_ser_trans1}, to
\[ n \leq {3m \over 2} \qquad \text{and} \qquad n \geq -{3m \over 2} \]
respectively. 
(For both summations we may take the lower limit to be zero since terms with $m=0$ or $n=0$ make no contribution to the sum.) 
It follows that
\[ S = \sum_{m=0}^\infty \sum_{n=-\lfloor {3m \over 2} \rfloor}^{\lfloor {3m \over 2} \rfloor} \left( {-4 \over 2m-n} \right) \left( {3m-2n \over 6} \right) \varphi(2m-n,3m-2n). \]
It is straightforward to check that the product of Kronecker symbols in the summand,
\[ \left( {-4 \over 2m-n} \right) \left( {3m-2n \over 6} \right), \]
is an odd function of $n$. 
The result of pairing terms of index $n$ in the inner sum with those of index $-n$ is therefore
\begin{equation} \label{eqn:S_rewritten}
S = \sum_{m=1}^\infty \sum_{n=1}^{\lfloor {3m \over 2} \rfloor} \left( {-4 \over 2m-n} \right) \left( {3m-2n \over 6} \right) \Big( \varphi(2m-n,3m-2n) - \varphi(2m+n,3m+2n) \Big).
\end{equation}
Now make the substitution~\eqref{eqn:unimodular_sub} once again. 
The inequalities
\[ n \geq 0 \qquad \text{and} \qquad n \leq {3m \over 2}, \]
which define the range of the summation in~\eqref{eqn:S_rewritten}, become
\[ n \leq {3m \over 2} \qquad \text{and} \qquad n \geq 0 \]
respectively. 
Hence, the right-hand side of~\eqref{eqn:S_rewritten} is identical to that of~\eqref{eqn:double_ser_trans1}.

Now, in order to establish the identity~\eqref{eqn:double_ser_trans2}, let
\[ T= \sum_{m=1}^\infty \sum_{n=3m+1}^{\lfloor {24 \over 7}m \rfloor} \left({-4 \over m}\right) \left({n \over 6}\right) \varphi(m,n). \]
The unimodular substitution $(m,n) \mapsto (7m-2n,24m-7n)$ changes the inequalities
\[ n > 3m \qquad \text{and} \qquad n \leq {24m \over 7} \]
to
\[ n < 3m \qquad \text{and} \qquad n \geq 0 \]
respectively. 
The result of applying this substitution to the series $T$ is therefore
\[ T = \sum_{m=1}^\infty \sum_{n=0}^{3m-1} \left({-4 \over 7m-2n}\right) \left({24m-7n \over 6}\right) \varphi(7m-2n,24m-7n). \]
It is straightforward to verify that, for any integers $m$ and $n$,
\[ \left({-4 \over 7m-2n}\right) \left({24m-7n \over 6}\right) = -\left({-4 \over m}\right) \left({n \over 6}\right). \]
Consequently,
\begin{equation} \label{eqn:T_rewritten}
T = -\sum_{m=1}^\infty \sum_{n=1}^{3m} \left({-4 \over m}\right) \left({n \over 6}\right) \varphi(7m-2n,24m-7n),
\end{equation}
the change of limits in the inner summation being justified since $({n\over6})$ is zero when $n=0$ or $3m$. 
The identity~\eqref{eqn:double_ser_trans2} now follows from adding the series
\[ \sum_{m=1}^\infty \sum_{n=1}^{3m} \left({-4 \over m}\right) \left({n \over 6}\right) \varphi(m,n) \]
to both sides of~\eqref{eqn:T_rewritten}.
\end{proof}

\end{document}
